% -----------------------------------------------------------------------------
% R1 exclusion: axisymmetric bounded-circulation residual class
% -----------------------------------------------------------------------------
% This file is written as a self-contained replacement/insert section.  It uses
% only standard PDE terminology.  The class R_1 is the axisymmetric
% bounded-circulation residual class from the refined residual stratification.

\section{Exclusion of the axisymmetric bounded-circulation residual class}
\label{sec:R1-exclusion}

In this section we prove the rigidity statement needed for the first refined
residual class.  The argument uses the scale-invariant angular circulation in
centered variables.  The essential point is that the centered Type~I equation
contains the Ornstein--Uhlenbeck drift.  Thus the angular-circulation equation
has an absorbing boundary at the symmetry axis together with a confining drift
at spatial infinity.  This gives a scalar Liouville theorem for bounded ancient
solutions of the circulation equation.

Throughout the section we write cylindrical coordinates as
\((\rho,\theta,Z)\), and an axisymmetric vector field as
\[
   V=V_\rho e_\rho+V_\theta e_\theta+V_Z e_Z .
\]
The centered Navier--Stokes equation is
\begin{equation}\label{eq:centered-R1}
   \partial_\tau V-\Delta V+\frac12 y\cdot\nabla V+\frac12 V
   +(V\cdot\nabla)V+\nabla P=0,
   \qquad \nabla\cdot V=0 .
\end{equation}
For an axisymmetric solution define the centered angular circulation
\begin{equation}\label{eq:centered-Gamma-def}
   \Gamma(\rho,Z,\tau):=\rho V_\theta(\rho,Z,\tau).
\end{equation}
This is the same scale-invariant quantity as the physical circulation
\(rU_\theta\).  If the physical representative is
\[
   U(x,t)=(-t)^{-1/2}V(x/\sqrt{-t},-\log(-t)),
\]
then
\[
   rU_\theta(r,z,t)=\rho V_\theta(\rho,Z,\tau),
   \qquad r=\sqrt{-t}\,\rho,
   \quad z=\sqrt{-t}\,Z.
\]

\begin{lemma}[Centered circulation equation]\label{lem:centered-circulation-equation}
Let \((V,P)\) be a smooth bounded axisymmetric solution of
\eqref{eq:centered-R1}.  Then \(\Gamma=\rho V_\theta\) satisfies
\begin{equation}\label{eq:centered-Gamma-equation}
   \partial_\tau\Gamma
   =
   \partial_\rho^2\Gamma-\frac1\rho\partial_\rho\Gamma
      +\partial_Z^2\Gamma
   -\left(V_\rho+\frac\rho2\right)\partial_\rho\Gamma
   -\left(V_Z+\frac Z2\right)\partial_Z\Gamma
\end{equation}
on \((0,\infty)\times\mathbb R\times\mathbb R\).  Moreover
\begin{equation}\label{eq:axis-boundary-Gamma}
   \Gamma(0,Z,\tau)=0
   \qquad\hbox{for all }(Z,\tau)\in\mathbb R\times\mathbb R .
\end{equation}
\end{lemma}

\begin{proof}
The angular component of \eqref{eq:centered-R1} is
\[
   \partial_\tau V_\theta
   -\left(\Delta-\frac1{\rho^2}\right)V_\theta
   +\frac12(\rho\partial_\rho+Z\partial_Z)V_\theta
   +\frac12V_\theta
   +V_\rho\partial_\rho V_\theta+V_Z\partial_ZV_\theta
   +\frac{V_\rho V_\theta}{\rho}=0 .
\]
Multiplying by \(\rho\), using
\[
   \rho\left(\Delta-\frac1{\rho^2}\right)V_\theta
   =\partial_\rho^2\Gamma-\frac1\rho\partial_\rho\Gamma+
      \partial_Z^2\Gamma,
\]
\[
   \rho\left[\frac12(\rho\partial_\rho+Z\partial_Z)V_\theta
      +\frac12V_\theta\right]
   =\frac12(\rho\partial_\rho+Z\partial_Z)\Gamma,
\]
and
\[
   \rho\left(V_\rho\partial_\rho V_\theta+V_Z\partial_ZV_\theta
      +\frac{V_\rho V_\theta}{\rho}\right)
   =V_\rho\partial_\rho\Gamma+V_Z\partial_Z\Gamma,
\]
gives \eqref{eq:centered-Gamma-equation}.  Finally, smoothness of a bounded
axisymmetric vector field at the axis implies \(V_\theta(\rho,Z,\tau)=O(\rho)\)
as \(\rho\downarrow0\), locally uniformly in \((Z,\tau)\).  Hence
\(\Gamma=\rho V_\theta=O(\rho^2)\), which gives
\eqref{eq:axis-boundary-Gamma}.
\end{proof}

We next record the scalar Liouville theorem used below.  It is stated in a
form that isolates precisely the PDE mechanism: bounded drift perturbations of
the centered circulation operator cannot sustain a bounded ancient solution
which vanishes on the symmetry axis.

\begin{lemma}[Axis-absorption Liouville theorem]
\label{lem:axis-absorption-liouville}
Let \(B=(B_\rho,B_Z)\) be a smooth vector field on
\((0,\infty)\times\mathbb R\times\mathbb R\) satisfying
\[
   \|B\|_{L^\infty}<\infty .
\]
Let \(G\in C^{2,1}((0,\infty)\times\mathbb R\times\mathbb R)\cap L^\infty\)
solve
\begin{equation}\label{eq:abstract-Gamma-equation}
   \partial_\tau G
   =
   \partial_\rho^2G-\frac1\rho\partial_\rho G+
      \partial_Z^2G
   -\left(B_\rho+\frac\rho2\right)\partial_\rho G
   -\left(B_Z+\frac Z2\right)\partial_ZG
\end{equation}
for all \(\tau\in\mathbb R\), and assume
\begin{equation}\label{eq:abstract-axis-zero}
   \lim_{\rho\downarrow0}G(\rho,Z,\tau)=0
   \qquad\hbox{locally uniformly in }(Z,\tau).
\end{equation}
Then \(G\equiv0\).
\end{lemma}

\begin{proof}
Set \(M_B=\|B\|_{L^\infty}\).  For \(R>1\) define
\begin{equation}\label{eq:radial-barrier}
   h_R(\rho)
   :=
   \frac{\displaystyle\int_0^{\min\{\rho,R\}} s
      \exp\left(\frac{s^2}{4}-M_Bs\right)\,ds}
   {\displaystyle\int_0^R s
      \exp\left(\frac{s^2}{4}-M_Bs\right)\,ds} .
\end{equation}
Then \(0\le h_R\le1\), \(h_R(0)=0\), \(h_R(R)=1\), and, on
\(0<\rho<R\),
\begin{equation}\label{eq:hR-ode}
   h_R''+\left(M_B-\frac\rho2-\frac1\rho\right)h_R'=0,
   \qquad h_R'>0 .
\end{equation}
Let
\[
   \mathcal L_\tau f
   :=
   \partial_\rho^2f-\frac1\rho\partial_\rho f+\partial_Z^2f
   -\left(B_\rho+\frac\rho2\right)\partial_\rho f
   -\left(B_Z+\frac Z2\right)\partial_Z f .
\]
Since \(B_\rho\ge -M_B\), \eqref{eq:hR-ode} gives
\begin{equation}\label{eq:hR-superharmonic}
   \mathcal L_\tau h_R
   =h_R''-\left(\frac1\rho+\frac\rho2+B_\rho\right)h_R'
   \le
   h_R''+\left(M_B-\frac\rho2-\frac1\rho\right)h_R'
   =0 .
\end{equation}
Thus \((\partial_\tau-\mathcal L_\tau)h_R\ge0\).

We now compare \(G\) to this barrier on the strip
\(\Omega_R=(0,R)\times\mathbb R\).  Let
\(M_G=\|G\|_{L^\infty}\).  We first record the killed-strip decay estimate used
by the comparison argument.  If \(q_{s,R}\) solves
\[
   \partial_\tau q_{s,R}=\mathcal L_\tau q_{s,R},
   \qquad q_{s,R}=0\hbox{ on }\rho=0,R,
   \qquad q_{s,R}(\rho,Z,s)=1,
\]
then, for every compact \(K\Subset\Omega_R\) and every fixed \(t\),
\begin{equation}\label{eq:killed-strip-decay}
   \lim_{s\to-\infty}\sup_{(\rho,Z)\in K}q_{s,R}(\rho,Z,t)=0 .
\end{equation}
To prove \eqref{eq:killed-strip-decay}, fix \(K\Subset\Omega_R\).  Choose
\(\delta>0\), \(Z_K>0\), and \(Z_1>Z_K+1\) so that
\[
   K\subset [\delta,R-\delta]\times[-Z_K,Z_K]
   \subset (0,R)\times(-Z_1,Z_1).
\]
On the finite cylinder
\[
   \Omega_{R,Z_1}:=(0,R)\times(-Z_1,Z_1)
\]
with Dirichlet data on the full parabolic boundary, the operator
\(\mathcal L_\tau\) is uniformly parabolic on
\([\delta/2,R-\delta/2]\times[-Z_1,Z_1]\), has bounded first-order coefficients
there, and has absorbing boundary at \(\rho=0,R\).  The boundary parabolic
Harnack inequality and the strong maximum principle give a time \(T_R>0\) and a
number \(\vartheta_R\in(0,1)\), depending only on
\(R,M_B,\delta,Z_K\) and not on \(Z_1\), such that every nonnegative solution
with initial data bounded by \(1\) and zero lateral data satisfies
\[
   0\le q(\cdot,s+T_R)\le 1-\vartheta_R
   \qquad\hbox{on }[\delta,R-\delta]\times[-Z_K,Z_K].
\]
Applying the same estimate after each time translation and using the maximum
principle between the compact set and the boundary gives
\[
   \sup_K q_{s,R}(\cdot,s+mT_R)
   \le (1-\vartheta_R)^m,\qquad m=1,2,\dots .
\]
For a fixed \(t>s\), choose \(m\) with
\(s+mT_R\le t<s+(m+1)T_R\).  The maximum principle and interior Harnack estimate
on the final interval \([s+mT_R,t]\) preserve the preceding bound up to a
constant depending only on \(R,M_B,\delta,Z_K,T_R\).  Letting \(s\to-\infty\)
forces \(m\to\infty\), proving decay on \(K\) for the truncated problem.  The
solutions on \(\Omega_{R,Z_1}\) increase monotonically to the killed-strip
solution as \(Z_1\to\infty\); passing to this limit gives
\eqref{eq:killed-strip-decay}.

Fix \(R>1\), \(s<t\), and set
\[
   W^+_{R}:=\frac{G}{M_G}-h_R,
   \qquad
   W^-_{R}:=-\frac{G}{M_G}-h_R,
\]
with the convention that the conclusion is immediate if \(M_G=0\).
By \eqref{eq:abstract-Gamma-equation} and \eqref{eq:hR-superharmonic},
\[
   (\partial_\tau-\mathcal L_\tau)W_R^\pm\le0
   \qquad\hbox{in }\Omega_R\times(s,t).
\]
On \(\rho=0\), \eqref{eq:abstract-axis-zero} and \(h_R(0)=0\) give
\(W_R^\pm\le0\).  On \(\rho=R\), \(h_R(R)=1\) and \(|G|\le M_G\) give again
\(W_R^\pm\le0\).  The parabolic comparison principle with the killed-strip
kernel therefore yields, for \(\tau=t\),
\[
   (W_R^\pm)_+(\rho,Z,t)
   \le q_{s,R}(\rho,Z,t)
   \qquad ((\rho,Z)\in\Omega_R).
\]
Letting \(s\to-\infty\) and using \eqref{eq:killed-strip-decay} gives
\(W_R^\pm\le0\) on compact subsets of \(\Omega_R\).  Since the compact subset
is arbitrary,
\begin{equation}\label{eq:G-bounded-by-hR}
   |G(\rho,Z,t)|\le M_G h_R(\rho)
   \qquad (0<\rho<R, Z\in\mathbb R,
   \ t\in\mathbb R).
\end{equation}
Finally, for each fixed \(\rho<\infty\), the denominator in
\eqref{eq:radial-barrier} tends to infinity super-exponentially as
\(R\to\infty\), while the numerator remains fixed.  Hence
\(h_R(\rho)\to0\).  Letting \(R\to\infty\) in
\eqref{eq:G-bounded-by-hR} gives \(G(\rho,Z,t)=0\) for every
\(\rho>0\), \(Z\in\mathbb R\), and \(t\in\mathbb R\).  The boundary value at
\(\rho=0\) is already zero by assumption.  Therefore \(G\equiv0\).
\end{proof}

\begin{theorem}[Exclusion of the axisymmetric bounded-circulation residual class]
\label{thm:R1-exclusion}
Let \(\Cax(\mathcal S;I,J)\) be the axisymmetric bounded-circulation residual
class in the refined residual decomposition.  Then
\begin{equation}\label{eq:R1-empty}
   \Cax(\mathcal S;I,J)=\varnothing .
\end{equation}
\end{theorem}

\begin{proof}
Assume, for contradiction, that
\(V\in\Cax(\mathcal S;I,J)\).  By definition of \(\Cax\), the
associated physical ancient representative is axisymmetric and has bounded
angular circulation.  Equivalently, in centered variables,
\[
   \Gamma(\rho,Z,\tau)=\rho V_\theta(\rho,Z,\tau)
\]
is bounded on \((0,\infty)\times\mathbb R\times\mathbb R\).  Since \(V\) is a
Seregin ancient limit, it is smooth and bounded in centered variables.  Hence
\(B=(V_\rho,V_Z)\) is bounded.

By \Cref{lem:centered-circulation-equation}, \(\Gamma\) satisfies
\eqref{eq:centered-Gamma-equation} and vanishes at the axis.  Applying
\Cref{lem:axis-absorption-liouville} with \(G=\Gamma\) and
\(B=(V_\rho,V_Z)\) gives
\[
   \Gamma\equiv0 .
\]
Therefore \(V_\theta\equiv0\) for \(\rho>0\), and by smoothness also on the
axis.  Thus the associated ancient solution is axisymmetric without swirl.
But the axisymmetric no-swirl class is one of the structured alternatives
removed before the residual class is defined.  This contradicts
\(V\in\mathcal R(\mathcal S;I,J)\), and therefore no such \(V\) exists.
\end{proof}

\begin{remark}[Why this does not assert the full bounded-swirl Liouville theorem]
The proof uses the centered Type~I equation, specifically the
Ornstein--Uhlenbeck drift terms \(-\rho\partial_\rho/2\) and
\(-Z\partial_Z/2\) in the scalar equation for \(\Gamma\).  It is therefore a
Type~I Seregin-limit argument.  It does not prove that every bounded ancient
axisymmetric solution of the unrenormalized Navier--Stokes equations with
bounded physical circulation is swirl-free.  The latter lacks the confining
centered drift and remains a different shell-transport problem.
\end{remark}
